\section{基于 LLM 的智能体}

\subsection{决策即语言建模}

LLM 时代的核心观察：\textbf{决策制定可以被重新表述为因果语言建模问题}。

\begin{figure}[H]
\centering
\includegraphics[width=0.95\textwidth]{slides-images/slide-050.jpg}
\caption{决策即生成式轨迹建模：轨迹概率可分解为环境转移动力学和智能体策略两部分。右上：Decision Transformer 的架构示意\protect\footnotemark}
\end{figure}
\footnotetext{Slides 第51页。Source: Chen et al. 2021。}

轨迹概率的分解：

$$p(\tau \mid g) = p(s_1, a_1, s_2, a_2, \ldots \mid g) = \prod_t \underbrace{p(s_t \mid s_{t-1}, a_t)}_{\text{环境转移动力学}} \times \underbrace{\pi(a_t \mid \tau_{\leq t}, g)}_{\text{智能体策略}}$$

其中：
\begin{itemize}
    \item $\tau$：完整轨迹（状态-动作序列）
    \item $g$：自然语言目标
    \item $s_t$：$t$ 时刻的环境状态
    \item $a_t$：$t$ 时刻的动作
\end{itemize}

关键洞察：智能体策略 $\pi(a_t \mid \tau_{\leq t}, g)$ 本质上就是一个\textbf{条件语言模型}——给定历史轨迹和目标，预测下一个动作。因此，可以直接用自回归 Transformer 来建模。

\subsection{提示循环：最简单的 LLM 智能体}

最简单的 LLM 智能体就是在循环中进行提示：

\begin{enumerate}
    \item 在提示中描述\textbf{动作空间}（如：可以点击、输入、移动鼠标）
    \item 提供\textbf{指令}
    \item 将\textbf{历史动作和观测}追加到提示中
    \item 让模型预测下一个动作
    \item 执行动作，获取新观测，回到步骤3
\end{enumerate}

\begin{figure}[H]
\centering
\includegraphics[width=0.95\textwidth]{slides-images/slide-052.jpg}
\caption{提示循环智能体的基本结构：指定动作空间、指令和历史后，模型预测下一个动作\protect\footnotemark}
\end{figure}
\footnotetext{Slides 第53页。}

\begin{knowledgebox}{提示循环的本质}
提示循环智能体本质上就是\textbf{Chain-of-Thought prompting 在循环中运行}。每一步的``思考''对应推理，``动作''对应与环境的交互。这建立了推理与智能体之间的直接联系——推理能力越强的模型，理论上也能成为越好的智能体。
\end{knowledgebox}

\subsection{本章小结}

LLM 将智能体问题简化为条件语言建模：决策可以表述为根据目标和历史预测下一个动作。最简单的实现是提示循环，本质上是 CoT 在交互式环境中的应用。

%% ========== Part 8: 评估基准 ==========

